\documentclass[10pt,a4paper]{amsart}
\usepackage[margin=19mm]{geometry}
\usepackage{amsmath,amssymb,mathtools}
\usepackage[scr=rsfs,cal=euler]{mathalfa}
\usepackage{xfrac}
\usepackage[unicode,hidelinks,bookmarks=false]{hyperref}
\newcommand{\ie}{i.e.\ }
\newcommand{\cf}{cf.\ }
\newcommand{\resp}{respectively\ }
\newcommand{\Subsection}{\subsection}
\providecommand{\nobreakdash}{\nobreak\hbox{-}}
\setlength{\emergencystretch}{2em}
\multlinegap0pt
\usepackage{tikz}
\usetikzlibrary{cd}
\usepackage{bm}
\usepackage{mathtools}
\usepackage{dsfont}
\usepackage{xcolor}
\usepackage{amssymb}
\usepackage{enumerate}
\usepackage{float}
\newtheorem{lemma}{Lemma}
\newtheorem{proposition}{Proposition}
\newcommand{\R}{\mathbb R}
\title{On Pseudo-Effectivity and Volumes of Adjoint Classes in K\"ahler Families with Projective Central Fiber}
\author{Christopher D. Hacon, Yi Li, Sheng Rao}
\date{Source: arXiv:2602.04158v2 (CC BY 4.0)}
\begin{document}
\maketitle

\noindent\emph{Source and licence.} On Pseudo-Effectivity and Volumes of Adjoint Classes in K\"ahler Families with Projective Central Fiber. Version arXiv:2602.04158v2 (CC BY 4.0), distributed by arXiv under the Creative Commons Attribution 4.0 International licence, http://creativecommons.org/licenses/by/4.0/, as recorded in the arXiv licence metadata for this exact version and shown on the abstract page https://arxiv.org/abs/2602.04158v2. Full licence notice: \texttt{licenses/arxiv-2602.04158-CC-BY-4.0.txt}.\par\smallskip

\noindent\textit{Editorial scope.} Counted unit: the article's proposition DefAmple (source0.tex lines 1072--1082). The article's complete original proof of it is delivered with it (source0.tex lines 1083--1146, one \texttt{proof} environment of 64 lines). No mathematics is written, repaired or re-derived: the statement and the proof are the article's own lines, in the article's own notation, and no retained line is substituted or re-flowed.  Retained uncounted as the source-local context the proof needs: lines 1053--1069.  Nothing was omitted from any retained span, so the package carries no omission record.  No retained line contains a citation, so the wrapper carries no bibliography and no bibliography entry.  The retained lines contain the article's own diagram code, which the wrapper typesets with the standard tikz packages; no image file is delivered and the package declares no asset dependency.  Editorial formatting only, in addition: an AMS \texttt{amsart} preamble that restates the article's own declarations and macros used by the retained lines, namely the environments \texttt{lemma}, \texttt{proposition} on separate counters, together with the standard packages those lines need.  The retained environments print in this document as Lemma 1, Proposition 1 in order of appearance, whereas the article prints the same items with the numbering of its own full document.  The complete native source of the article as distributed in the arXiv e-print for this exact version is bundled unchanged as \texttt{sources/193958/source0.tex}.
\begin{lemma}\label{ResEmbed}
Let $p:X \to S$ and $g:Z \to  Y$ be proper morphisms of complex analytic varieties where $p$ is flat, and let $\Phi: X\to Z$ be a morphism such that the diagram commutes 
\begin{center}
\begin{tikzcd}[ampersand replacement=\&]
	X \&\& Z \\
	\& Y \\
	\& S.
	\arrow["\Phi", from=1-1, to=1-3]
	\arrow["f", from=1-1, to=2-2]
	\arrow["p"', from=1-1, to=3-2]
	\arrow["g"', from=1-3, to=2-2]
	\arrow["q", from=1-3, to=3-2]
	\arrow["r"{description}, from=2-2, to=3-2]
\end{tikzcd}
\end{center}
Fix a point $0\in S$ and assume that the restriction of the diagram $\Phi_0: X_0 \rightarrow Z_0/Y_0$ over $0\in S$ is a closed immersion (resp. an isomorphism). Then $\Phi:X\to Z /Y$ is a closed immersion (resp. an isomorphism) over some open neighborhood $U\subset S$ of $0$. 
\end{lemma}

\begin{proposition}\label{DefAmple}
Let $f: X \to Y/S$ be a proper surjective contraction morphism of complex analytic varieties over a smooth connected curve $S$. For each $t \in S$, denote by
\[
f_t : X_t \to Y_t
\]
the induced proper contraction on the fiber over $t$. 

(1) If $L$ is a $\mathbb{Q}$-line bundle on $X$ such that $L|_{X_0}$ is $f_0$-ample, then there exists a Zariski open subset $U \subset S$ around $0$ such that $L|_{X_t}$ is $f_t$-ample for any $t \in U$.

(2) If $\alpha\in H^{1,1}_{\rm{BC}}(X, \R)$ is a class on $X$ such that $\alpha|_{X_0}$ is $f_0$-K\"ahler, then there exists a (Euclidean) open subset $U\subset S$ around $0$ such that $\alpha|_{X_t}$ is $f_t$-K\"ahler for any $t\in U$. 
\end{proposition}

\begin{proof}
Since $L_0 = L|_{X_0}$ is $f_0$-ample,  there exists some $m\ge 0$ such that
$$\rho_m : f^* f_* (L^{\otimes m})\to L^{\otimes m}$$is surjective near $X_0$. To see this, consider the exact sequence 
\[
0 \to L^{\otimes m} \otimes \mathcal{I}_{X_0} \to L^{\otimes m} \to L^{\otimes m}|_{X_0} \to 0.
\]
Since $L|_{X_0}$ is $f_0$-ample, Serre’s vanishing theorem gives 
\[
R^i (f_0)_* \bigl(L^{\otimes m}|_{X_0}\bigr) = 0 \quad \text{for } i>0 \text{ and } m \gg 0.
\]
Therefore,
\[
\psi: R^1 f_* (L^{\otimes m}) \otimes \mathcal{I}_{\{0\}} 
\cong R^1 f_*\bigl(L^{\otimes m} \otimes \mathcal{I}_{X_0}\bigr) 
\longrightarrow R^1 f_* (L^{\otimes m})
\]
is surjective for $m \gg 0$, where the first isomorphism follows from the projection formula together with the fact that $\mathcal{I}_{X_0}=f^*\mathcal{I}_{\{0\}}$ and $\mathcal{I}_{\{0\}}$ are invertible on the smooth connected curve $S$. Taking the stalk at $0$ and tensor with $\kappa(0)\cong \mathbb{C}$, we obtain a surjective map
\[
\psi(0): \;  R^1 f_* (L^{\otimes m})_0 \otimes \bigl(\mathfrak{m}_0 / \mathfrak{m}_0^2\bigr)
\longrightarrow R^1 f_* (L^{\otimes m})_0 \otimes \kappa(0).
\]
Note that this is the zero map, since the composition in the short exact sequence 
\[
0\to \mathfrak{m}_{0} /\mathfrak{m}_{0}^2 \to \mathcal{O}_{T,0} /\mathfrak{m}_{0}^2 \to \mathcal{O}_{T,0} /\mathfrak{m}_{0} = \kappa (0) \to 0
\] 
is zero. Therefore, $R^1 f_* (L^{\otimes m})_0 \otimes \kappa(0) = 0$ as well. By Nakayama’s lemma, we conclude that 
\[
R^1 f_* (L^{\otimes m})_0 = 0.
\]
Since $R^1 f_* (L^{\otimes m})$ is coherent, its support is a closed analytic subset. Therefore, $R^1 f_* (L^{\otimes m}) = 0$ in a neighborhood of $0$. Consequently,
\[
f_* L^{\otimes m} \longrightarrow (f_0)_* \bigl(L^{\otimes m}|_{X_0}\bigr)
\]
is surjective. Since $L^{\otimes m}|_{X_0}$ is $f_0$-globally generated for $m\gg 0$, i.e., the map
\[
f_0^* (f_0)_* (L^{\otimes m} |_{X_0} ) \to L^{\otimes m}|_{X_0}
\]
is surjective for $m\gg 0$, the morphism
\[
\rho_m: f^* f_* L^{\otimes m} \longrightarrow L^{\otimes m}, \quad \text{for } m\gg 0
\]
is surjective in a neighborhood of $X_0$.


Therefore, after shrinking $S$, it will induce the commutative diagram  
\begin{center}
\begin{tikzcd}[ampersand replacement=\&]
	X \&\& {\mathbb{P}_Y(f_* L^{\otimes m})} \\
	\& Y \\
	\& S.
	\arrow["\Phi", from=1-1, to=1-3]
	\arrow["f"{description}, from=1-1, to=2-2]
	\arrow["g"', from=1-1, to=3-2]
	\arrow[from=1-3, to=2-2]
	\arrow[from=1-3, to=3-2]
	\arrow["h"{description}, from=2-2, to=3-2]
\end{tikzcd}
\end{center}
Since $L|_{X_0}$ is $f_0$-ample, the restriction to the fiber over $0 \in S$ induces an embedding
$$\Phi_0 = \Phi|_{X_0}: X_0 \hookrightarrow \mathbb{P}_{Y_0}((f_0)_* (L^{\otimes m}|_{X_0}))/Y_0.$$
By Lemma \ref{ResEmbed}, after further shrinking the base $S$, $\Phi:X \hookrightarrow \mathbb{P}_{Y}(f_* L^{\otimes m})/Y$ is still a closed embedding.

The second part of the proposition is clear, since positive form is still positive when restricted to the nearby fibers. 
\end{proof}

\end{document}
